\DocumentMetadata{
  pdfversion=2.0,
  pdfstandard=ua-2,
  lang=en-US,
  tagging=on,
  tagging-setup={math/setup=mathml-SE,tabsorder=structure}
}
\documentclass[11pt,letterpaper]{article}
\usepackage[margin=0.68in,includefoot]{geometry}
\usepackage{newcomputermodern}
\usepackage{amsmath,amscd,mathtools}
\usepackage{tikz,adjustbox}
\providecommand{\square}{\mdlgwhtsquare}
\usepackage[hidelinks]{hyperref}
\hypersetup{
  pdftitle={Algebra PhD Exam, August 2020},
  pdfauthor={University of Florida Department of Mathematics},
  pdfsubject={Graduate mathematics examination},
  pdfdisplaydoctitle=true,
  bookmarksopen=true
}
\setcounter{secnumdepth}{0}
\setlength{\parindent}{0pt}
\setlength{\parskip}{0.28em}
\setlength{\emergencystretch}{3em}
\setlength{\leftmargini}{2.15em}
\setlength{\leftmarginii}{2.65em}
\setlength{\leftmarginiii}{3em}
\pagestyle{empty}
\raggedbottom
\allowdisplaybreaks
\NewTaggingSocketPlug{block/list/label}{examproblem}{%
  \tagstructbegin{tag=Lbl}\tagstructend%
  \tagstructbegin{tag=\LBody}%
  \tagstructbegin{tag=\ProblemHeadingTag,title-o={\ProblemHeadingText}}%
  \tagmcbegin{tag=\ProblemHeadingTag,actualtext={\ProblemHeadingText}}%
  #2%
  \tagmcend%
  \tagstructend%
}
\newenvironment{examproblems}[2]{%
  \def\ProblemHeadingTag{#2}%
  \AssignTaggingSocketPlug{block/list/label}{examproblem}%
  \begin{enumerate}%
  \setcounter{enumi}{#1}%
  \renewcommand{\labelenumi}{\textbf{\theenumi.}}%
  \setlength{\topsep}{0.35em}%
  \setlength{\partopsep}{0pt}%
  \setlength{\itemsep}{0.48em}%
  \setlength{\parsep}{0.18em}%
}{\end{enumerate}}
\newenvironment{examsubparts}{%
  \AssignTaggingSocketPlug{block/list/label}{default}%
  \begin{enumerate}%
  \setlength{\topsep}{0.18em}%
  \setlength{\partopsep}{0pt}%
  \setlength{\itemsep}{0.16em}%
  \setlength{\parsep}{0.1em}%
}{\end{enumerate}}
\newenvironment{examinstructions}{%
  \par\begingroup\itshape
}{%
  \par\endgroup\vspace{0.18em}
}
\makeatletter
\renewcommand\subsection{\@startsection{subsection}{2}{0pt}%
  {-1.25ex plus -.25ex minus -.1ex}%
  {0.55ex plus .1ex}%
  {\normalfont\large\bfseries}}
\renewcommand{\@maketitle}{%
  \newpage\null\vskip -1em%
  {\footnotesize\bfseries
    \href{https://www.ufl.edu/}{University of Florida}, \href{https://math.ufl.edu/}{Department of Mathematics}\par}%
  \vskip -0.5em%
  \tagstructbegin{tag=H1,title={Algebra PhD Exam, August 2020}}%
  \tagpdfparaOff%
  \tagmcbegin{tag=H1}%
    {\LARGE\bfseries\href{https://gma.math.ufl.edu/exams/algebra-phd/}{Algebra PhD Exam}, August 2020\par}%
  \tagmcend%
  \tagpdfparaOn%
  \tagstructend%
  \vskip 0.25em%
  \tagmcbegin{artifact}%
    \hrule height 1pt%
  \tagmcend%
  \vskip 1em%
}
\makeatother
\title{Algebra PhD Exam, August 2020}
\author{}
\date{}
\begin{document}
\maketitle
\thispagestyle{empty}
\begin{examinstructions}
Answer seven problems. (If you turn in more, the first seven will be graded.) Within reason, you may quote theorems as long as you state them clearly.
\end{examinstructions}
\begin{examinstructions}
Below ring means associative ring with identity, and module means unital module unless otherwise specified.
\end{examinstructions}
\begin{examproblems}{0}{H2}
\def\ProblemHeadingText{Problem 1.}
\item
Let~\(K\) be a finite field, and let~\(F\) be a subfield of~\(K\). Prove that the extension \(K/F\) is Galois.
\def\ProblemHeadingText{Problem 2.}
\item
Let \(p(x)=x^7-3\in\mathbb{Q}[x]\), \(\zeta\in\mathbb{C}\) a primitive 7-th root of 1, \(\xi=\sqrt[7]{3}\in\mathbb{C}\), and~\(K\) the splitting field of \(p(x)\).
\begin{examsubparts}
\item[\textnormal{(a)}]
Compute the Galois group of \(K/\mathbb{Q}\).
\item[\textnormal{(b)}]
Compute the Galois group of \(K(\zeta)/\mathbb{Q}(\zeta)\).
\item[\textnormal{(c)}]
Compute the Galois group of \(K(\xi)/\mathbb{Q}(\xi)\).
\end{examsubparts}
\end{examproblems}
\begin{examinstructions}
In every case, carefully justify that your calculation is correct.
\end{examinstructions}
\begin{examproblems}{2}{H2}
\def\ProblemHeadingText{Problem 3.}
\item
Let~\(A\) and~\(B\) be abelian groups, and \(T=A\otimes_{\mathbb{Z}}B\). Let \(\phi:A\to A\) be a group homomorphism. Prove that there is a group homomorphism
\[
\psi:T\to T
\]
 such that, for all \(a\in A\) and \(b\in B\), we have
\[
\psi(a\otimes_{\mathbb{Z}}b)=\phi(a)\otimes_{\mathbb{Z}}b.
\]
\def\ProblemHeadingText{Problem 4.}
\item
Let \(\mathcal{FG}\) be the category of all finite groups.
\begin{examsubparts}
\item[\textnormal{(a)}]
Recall the definition of free object in an arbitrary concrete category~\(\mathcal{C}\).
\item[\textnormal{(b)}]
Prove from your definition that there exists a finite set~\(S\) such that there is no free object in \(\mathcal{FG}\) freely generated by~\(S\).
\end{examsubparts}
\def\ProblemHeadingText{Problem 5.}
\item
Give an example of a projective module which is not free. Prove that your example has the desired properties.
\def\ProblemHeadingText{Problem 6.}
\item
State and prove Hilbert’s Basis Theorem.
\def\ProblemHeadingText{Problem 7.}
\item
Prove the Lying-Over Theorem: Let~\(S\) be an integral extension of an integral domain~\(R\), and let~\(P\) be a prime ideal of~\(R\). Then, there exists a prime ideal~\(Q\) of~\(S\), such that \(Q\cap R=P\).
\def\ProblemHeadingText{Problem 8.}
\item
Let~\(R\) be an integral domain that is a local ring. Let~\(I\) be an invertible ideal of~\(R\). Prove that~\(I\) is principal.
\def\ProblemHeadingText{Problem 9.}
\item
Let~\(R\) be the ring of all rational numbers with odd denominators. Prove that the Jacobson radical \(J(R)\) of~\(R\) consists of all fractions whose denominator is odd and whose numerator is even.
\def\ProblemHeadingText{Problem 10.}
\item
This problem has three parts.
\begin{examsubparts}
\item[\textnormal{(a)}]
Define what is a Dedekind domain.
\item[\textnormal{(b)}]
Give an example of a Dedekind domain which is not a unique factorization domain.
\item[\textnormal{(c)}]
Prove that your example is a Dedekind domain and not a unique factorization domain carefully explaining any standard results you use.
\end{examsubparts}
\def\ProblemHeadingText{Problem 11.}
\item
Let~\(R\) be a ring with identity. Consider the ring \(\operatorname{Hom}_R(R,R)\) of left~\(R\)-module homomorphisms from~\(R\) to itself. Prove that \(\operatorname{Hom}_R(R,R)\) is isomorphic, as a ring, to \(R^{\mathrm{op}}\).
\end{examproblems}
\end{document}
